%% ma: L12-B07-0022
%% lop: 12
%% chuong: 2
%% bai: 7
%% dang: 12.7.6
%% loai: TN4
%% muc: NB
%% dap_an: "C"
%% nguon: "CĐ2.2.pdf giáo viên gửi 27/10/2026 (M07, Lê Minh Khôi – Nguyễn Anh Khôi, Chương II Chủ đề 2 Tọa độ của vectơ trong không gian), Câu 19"
%% kien_thuc: "Bài 7 (hình chiếu của điểm lên mặt phẳng tọa độ)"
%% trang_thai: chinh-thuc
\begin{ex}
Trong không gian $Oxyz$, cho điểm $P(2026;2027;10)$. Gọi $P'$ là điểm đối xứng của $P$ qua mặt phẳng $(Oxz)$. Tọa độ của điểm $P'$ là
\choice{$P'(-2026;-2027;10)$}{$P'(2026;2027;-10)$}{\True $P'(2026;-2027;10)$}{$P'(0;2027;0)$}
\loigiai{Đối xứng qua $(Oxz)$: đổi dấu tung độ: $P'(2026;-2027;10)$.}
\end{ex}
